\chapter{Compensation des swaps et base entre chambres de compensation}\label{ch:m2:swap-clearing-and-the-ccp-basis}

En 2009, les dirigeants du G20 sont convenus que tout dérivé de gré à gré standardisé
devrait être compensé par une contrepartie centrale d'ici à la fin de
2012. Les swaps, qui étaient des promesses bilatérales entre banques, sont devenus des transactions
compensées par une chambre de compensation, avec appels de marge quotidiens, et le prix d'un swap était censé
ne dépendre que de ses modalités. Il n'en a pas tout à fait été ainsi. En 2017, un swap
de trente ans en dollars, compensé dans une chambre de compensation, était coté avec plusieurs
points de base d'écart par rapport au même swap compensé dans l'autre, et une société de données
pouvait montrer que, pour 100~millions de dollars, l'écart représentait environ 727
milliers de dollars de financement de marge. Les flux de trésorerie du swap sont les mêmes~; ce qui
diffère, c'est l'identité des autres intervenants de chaque chambre, et donc la marge qu'un
dealer doit déposer pour tenir l'autre côté de la transaction. Ce chapitre explique la compensation,
la marge qui l'accompagne, et pourquoi un même swap peut avoir deux prix.

\section{Obligation de compensation}

\begin{definition}[Obligation de compensation, compensation pour compte de clients et portage]\label{def:m2:swap-clearing-and-the-ccp-basis:mandate}
Une \emph{obligation de compensation}\index{obligation de compensation} est une exigence réglementaire
qui impose que certaines catégories désignées de dérivés soient compensées par une
contrepartie centrale (livre~1, chapitre~5). Les entreprises qui ne sont pas membres de la
chambre de compensation passent par un membre compensateur, qui garantit leurs transactions~: c'est
\emph{compensation pour compte de clients}\index{compensation pour compte de clients}. Le \emph{portage des positions}\index{portage des positions} est
le transfert des positions et de la marge d'un client, d'un membre compensateur défaillant
à un autre, afin que les couvertures du client survivent à la défaillance de son courtier.
\end{definition}

Les États-Unis ont été les premiers à mettre en œuvre cet engagement~: la Commodity Futures
Trading Commission a adopté sa première exigence de compensation en novembre 2012, pour
quatre catégories de \omterm{def:m2:interest-rate-swaps:swap}{swaps de taux} d'intérêt et deux de swaps sur indices de crédit, et les
dealers et les fonds actifs ont commencé à les compenser en mars 2013. L'Union européenne a suivi
avec sa propre réglementation et a ajouté en 2024 une exigence selon laquelle ses entreprises
doivent détenir un compte actif auprès d'une chambre de compensation située dans l'Union, un choix
politique autant qu'une règle de risque, quant au lieu où les swaps en euros sont compensés.

Ce que change la compensation, c'est la forme du réseau
(\cref{fig:m2:swap-clearing-and-the-ccp-basis:network}). Cinq dealers qui négocient de gré à gré
font face à dix relations de contrepartie, chacune avec son propre accord de collatéral
et sa propre exposition~; compensés, chacun ne fait face qu'à une chambre de compensation, dont
les transactions avec tous les autres se nettent en une seule position, et un défaut est
absorbé par la marge et le fonds de défaillance de la chambre plutôt que transmis le long de la
toile.

\begin{omfigure}
\centering
\begin{tikzpicture}[font=\small,
  dl/.style={circle,draw=omDef,thick,fill=omDef!10,minimum size=7mm,inner sep=0pt}]
\foreach \i in {1,...,5} {\node[dl] (a\i) at ({90+72*(\i-1)}:1.6) {D\i};}
\foreach \i/\j in {1/2,1/3,1/4,1/5,2/3,2/4,2/5,3/4,3/5,4/5} {\draw[gray,thick] (a\i) -- (a\j);}
\node[font=\footnotesize,align=center] at (0,-2.5) {bilatérales~: 10 relations};
\begin{scope}[xshift=6.5cm]
\node[rectangle,draw=omThm,thick,fill=omThm!12,minimum size=9mm,rounded corners=2pt] (c) at (0,0) {CCP};
\foreach \i in {1,...,5} {\node[dl] (b\i) at ({90+72*(\i-1)}:1.6) {D\i}; \draw[omThm,thick] (b\i) -- (c);}
\node[font=\footnotesize,align=center] at (0,-2.5) {compensées~: 5 relations};
\end{scope}
\end{tikzpicture}

\omcaption{Cinq dealers négociant des swaps de gré à gré, chaque paire ayant son propre
accord de collatéral, et les mêmes dealers après compensation, chacun ne faisant face qu'à la
chambre de compensation. Schéma de principe.}\label{fig:m2:swap-clearing-and-the-ccp-basis:network}
\end{omfigure}

\begin{dated}{2026-09}{Obligations de compensation et règles de marge}\label{dat:m2:swap-clearing-and-the-ccp-basis:mandates}
\textbf{États-Unis}~: compensation obligatoire pour les catégories de swaps désignées à partir du
11 mars 2013 pour les dealers et les fonds actifs, puis pour les autres. \textbf{Union européenne}~:
en vertu d'EMIR 3, en vigueur depuis le 24 décembre 2024, les entreprises soumises à l'\omterm{def:m2:swap-clearing-and-the-ccp-basis:mandate}{obligation
de compensation} doivent détenir un compte actif auprès d'une chambre de compensation agréée
dans l'UE et y compenser un nombre représentatif de dérivés de taux d'intérêt en euros et en zlotys
par son intermédiaire, à partir du 24 juin 2025. \textbf{Transactions non compensées}~: le
cadre international de marges a atteint sa phase finale le 1\up{er} septembre 2022, pour les
entreprises détenant plus de 8~milliards d'euros de dérivés non compensés~; la marge initiale n'a
pas à être échangée en deçà de 50~millions d'euros par groupe de contreparties.
\end{dated}

\section{Les chambres de compensation en concurrence}

Les swaps en dollars sont compensés principalement dans deux chambres, LCH's SwapClear et CME~; les
autres devises ont leur propre répartition, et une partie de la compensation en euros doit désormais
avoir lieu au sein de l'Union (\cref{dat:m2:swap-clearing-and-the-ccp-basis:mandates}). Un dealer est membre de plusieurs chambres et peut choisir où compenser
une transaction avec un autre dealer. Un client compense généralement là où son courtier
compensateur et ses propres systèmes de risque sont en place, et il ne peut pas déplacer librement ses transactions entre
chambres~: la marge est calculée séparément sur le portefeuille de chaque chambre, et une
transaction dans l'une ne se compense jamais avec une transaction dans l'autre.

\begin{omfigure}
\begin{tikzpicture}[font=\small,
  box/.style={draw,thick,rounded corners=2pt,align=center,minimum height=1cm,minimum width=2.9cm,inner sep=4pt},
  arr/.style={-{Stealth[length=2.5mm]},thick}]
\node[box,draw=omDef,fill=omDef!8] (cl) at (0,0) {clients\\(surtout payeurs fixes)};
\node[box,draw=omThm,fill=omThm!8] (dA) at (5,1.3) {dealer auprès de la CCP A\\reçoit le taux fixe};
\node[box,draw=omThm,fill=omThm!8] (dB) at (5,-1.3) {dealer auprès de la CCP B\\paie le taux fixe};
\node[box,draw=gray,fill=gray!8] (st) at (10.2,-1.3) {autres dealers\\(reçoivent le taux fixe)};
\draw[arr] (cl) -- node[above left=2pt,align=center] {paie le taux fixe\\auprès de la CCP A} (dA);
\draw[arr,dashed] (dA) -- node[right=4pt,align=left] {même dealer,\\couverture} (dB);
\draw[arr] (dB) -- node[below=2pt,align=center] {paie le taux fixe\\auprès de la CCP B} (st);
\node[font=\footnotesize,text=omProp,align=center] at (10.2,1.3) {marge déposée deux fois~:\\une fois auprès de la CCP A, une fois auprès de la CCP B};
\end{tikzpicture}

\omcaption{Origine de la \omterm{def:m2:swap-clearing-and-the-ccp-basis:basis}{base entre chambres de compensation}. Les clients d'une chambre sont surtout
payeurs fixes~; les dealers qui prennent l'autre côté reçoivent le taux fixe là-bas et se
couvrent en payant le taux fixe dans l'autre chambre, où d'autres dealers reçoivent. Le
risque du dealer s'annule sur les deux, mais chaque chambre appelle la marge sur sa moitié
de son côté, et le dealer finance les deux marges.}\label{fig:m2:swap-clearing-and-the-ccp-basis:flows}
\end{omfigure}

\section{Marge initiale, compensée et non compensée}

Une chambre de compensation collecte la marge de variation, soit la variation quotidienne de la
valeur d'un swap, et la marge initiale, un coussin contre la perte qu'elle subirait en
liquidant le portefeuille d'un membre défaillant (livre~1, chapitre~20).

\begin{definition}[Règles de marge pour les dérivés non compensés et credit support annex]\label{def:m2:swap-clearing-and-the-ccp-basis:umr}
Les \emph{règles de marge pour les dérivés non compensés}\index{règles de marge pour les dérivés non compensés} sont les
réglementations qui obligent les grands utilisateurs de dérivés qui ne sont pas
compensés de manière centrale à échanger bilatéralement leurs marges initiales et de variation, la marge initiale
étant conservée par un tiers. Une \emph{credit support annex}\index{credit support annex}
(CSA) est la partie d'un accord bilatéral sur les dérivés qui régit le
collatéral~: quels actifs, quels seuils, quels transferts minimaux, le taux versé
sur le collatéral en espèces, et le calendrier des appels.
\end{definition}

\begin{proposition}[Marge initiale et son coût]\label{prop:m2:swap-clearing-and-the-ccp-basis:mva}
Supposons qu'une banque appelle la marge sur une position de swap à $z$ écarts types de la
variation de sa valeur sur une période de marge en risque de $h$ jours, avec une volatilité quotidienne des taux de
$\sigma$ points de base~: $\mathrm{IM}(t) = z\,\sigma\sqrt{h}\,\mathrm{DV01}(t)$. Si
la constitution de la marge coûte un spread de financement $s$ par an, l'\emph{ajustement de
valeur de marge} est
\[
\mathrm{MVA} \;=\; \int_0^T s\,\mathrm{IM}(t)\,P(0,t)\,dt,
\]
et le spread courant du swap qui le rémunère est
$\mathrm{MVA}/\mathrm{DV01}(0)$ points de base.
\end{proposition}
\begin{proof}
La première affirmation définit le modèle~: les variations de valeur sont normales, d'écart type
$\sigma\sqrt h$ points de base de taux, multiplié par le DV01. Le coût
de détenir $\mathrm{IM}(t)$ pendant $dt$ est $s\,\mathrm{IM}(t)\,dt$, actualisé. Un
spread courant de $b$ points de base sur le swap vaut $b\,\mathrm{DV01}(0)$
aujourd'hui.
\end{proof}

\begin{example}[Un swap à dix ans et un swap à trente ans]\label{ex:m2:swap-clearing-and-the-ccp-basis:mva}
Avec des paramètres illustratifs, $\sigma = 7$ points de base par jour, une période de marge de cinq jours
et $z = 2.326$, la marge initiale vaut 36.4 points de base de DV01. Sur 100~millions
de dollars d'un swap à dix ans à 4~\%, le DV01 vaut 81\,109~dollars, la marge initiale démarre à
2.95~millions de dollars et décroît avec la durée de vie restante du swap
(voir la \cref{fig:m2:swap-clearing-and-the-ccp-basis:im})~; financée à 50 points de base par an
elle coûte 76\,442~dollars sur la durée de vie du swap~: 0.94 point de base par an de
spread courant. Pour le swap à trente ans, 2.30 points de base.
\end{example}

\begin{omfigure}
\begin{tikzpicture}
\begin{axis}[width=11cm,height=4.8cm,
  xlabel={années depuis la transaction},ylabel={marge initiale (M USD)},
  tick label style={font=\small},label style={font=\small},
  xmin=0,xmax=30,ymin=0,ymax=7,grid=major,grid style={gray!25},
  legend columns=2,legend style={font=\footnotesize,at={(0.5,-0.3)},anchor=north,draw=none,fill=none,/tikz/every even column/.append style={column sep=8pt}},
  table/col sep=comma]
\addplot[omDef,very thick,restrict x to domain=0:9.9] table[x=t,y=ten]{figdata/markets-2/10-swap-clearing-and-the-ccp-basis/im.csv};
\addplot[omThm,very thick] table[x=t,y=thirty]{figdata/markets-2/10-swap-clearing-and-the-ccp-basis/im.csv};
\legend{swap à 10 ans,swap à 30 ans}
\end{axis}
\end{tikzpicture}

\omcaption{Marge initiale sur la durée de vie de 100~millions de dollars d'un swap à dix ans et
d'un swap à trente ans, dans le modèle stylisé de la
\cref{prop:m2:swap-clearing-and-the-ccp-basis:mva}. La marge suit le DV01 du
swap restant, et l'aire sous chaque courbe, multipliée par le spread de financement,
en donne le coût. Les paramètres sont illustratifs. Données~: le tutoriel du
chapitre.}\label{fig:m2:swap-clearing-and-the-ccp-basis:im}
\end{omfigure}

\section{Pourquoi un même swap a deux prix}

\begin{definition}[Base entre chambres de compensation]\label{def:m2:swap-clearing-and-the-ccp-basis:basis}
La \emph{base entre chambres de compensation}\index{base entre chambres de compensation} est l'écart entre les \omterm{def:m2:government-bonds:zero}{taux au pair}
de swaps économiquement identiques compensés dans deux chambres de compensation différentes.
\end{definition}

Si les transactions d'un dealer dans une chambre se compensent entre elles, il y poste peu
de marge initiale là-bas. Si ses clients de cette chambre paient tous le taux fixe, il y est un receveur unidirectionnel
là-bas, et ses couvertures, qui paient le taux fixe dans l'autre chambre, le laissent
payeur unidirectionnel là-bas~: deux marges, toutes deux financées, pour une position neutre. Le dealer
facture ce coût lorsqu'il reçoit le taux fixe dans la première chambre, en demandant
un taux fixe plus élevé à cet endroit. L'analyse de 2017 a montré que l'effet est le plus marqué pour les
swaps en dollars à long terme, là où les clients payeurs chez CME étaient plus nombreux que les receveurs~: environ 3.4 points de base
à trente ans, dont 1.3 correspondait au coût de marge chez CME et 2.1 chez LCH.

\begin{omfigure}
\begin{tikzpicture}
\begin{axis}[width=11cm,height=4.8cm,
  xlabel={échéance du swap (années)},ylabel={base par position unidirectionnelle (pb)},
  tick label style={font=\small},label style={font=\small},
  xmin=0,xmax=31,ymin=0,ymax=3.7,grid=major,grid style={gray!25},
  legend columns=3,legend style={font=\footnotesize,at={(0.5,-0.3)},anchor=north,draw=none,fill=none,/tikz/every even column/.append style={column sep=8pt}},
  table/col sep=comma]
\addplot[omDef,thick,dashed,mark=*,mark size=1.4pt] table[x=maturity,y=f25]{figdata/markets-2/10-swap-clearing-and-the-ccp-basis/basis.csv};
\addplot[omThm,very thick,mark=*,mark size=1.4pt] table[x=maturity,y=f50]{figdata/markets-2/10-swap-clearing-and-the-ccp-basis/basis.csv};
\addplot[omProp,thick,dotted,mark=*,mark size=1.4pt] table[x=maturity,y=f75]{figdata/markets-2/10-swap-clearing-and-the-ccp-basis/basis.csv};
\legend{financement 25 pb,financement 50 pb,financement 75 pb}
\end{axis}
\end{tikzpicture}

\omcaption{Le spread courant qui rémunère une position unidirectionnelle margée, par
échéance et par spread de financement, dans le modèle stylisé. Un dealer ayant des positions unidirectionnelles
dans deux chambres en a besoin deux fois plus. La base croît avec
l'échéance, car la marge est détenue plus longtemps, mais moins que proportionnellement,
car la marge s'amortit avec le swap et les années lointaines sont davantage actualisées. Données~:
le tutoriel du chapitre.}\label{fig:m2:swap-clearing-and-the-ccp-basis:basis}
\end{omfigure}

La base n'est donc pas un arbitrage. Recevoir le taux fixe dans la chambre où il est cher et
payer dans l'autre, où il est bon marché, verrouille le spread, mais crée aussi les deux
positions unidirectionnelles dont le coût de marge est précisément ce que le spread rémunère~; seule une firme dont les
positions existantes dans les deux chambres font baisser la marge exigée par les nouvelles
transactions peut le capter, et les dealers se font concurrence pour être cette firme. Les mandats qui orientent les flux vers une
chambre particulière, comme le compte actif de l'UE, déplacent eux aussi la base.

\section{Tutoriel~: valorisation de la marge}\label{tut:m2:swap-clearing-and-the-ccp-basis:mva}

\begin{tutorial}
\textbf{Objectif.} Calculer le profil de marge initiale d'un swap, son coût de financement
et la base qui le rémunère. \textbf{État final~:}
\cref{fig:m2:swap-clearing-and-the-ccp-basis:im,fig:m2:swap-clearing-and-the-ccp-basis:basis}
et les résultats de l'\cref{ex:m2:swap-clearing-and-the-ccp-basis:mva}.
\begin{tutsteps}
\item \textbf{Le modèle de marge}~: marge initiale par unité de DV01.
\omcode{code/firm/ccpbasis/firm_ccpbasis.py}{24}{32}{Modèle stylisé de marge initiale.}\label{lst:m2:swap-clearing-and-the-ccp-basis:model}
\item \textbf{Marge sur la durée de vie, son coût et la base.}
\omcode{code/firm/ccpbasis/firm_ccpbasis.py}{44}{58}{Profil de marge initiale, ajustement de valorisation de marge (MVA) et base courante.}\label{lst:m2:swap-clearing-and-the-ccp-basis:mva}
\item \textbf{Exécuter} \texttt{ccp\_demo.one\_way\_book()} et
\texttt{fig\_ccp.py}.
\end{tutsteps}
\textbf{Pour aller plus loin.} Doubler la période de marge en risque pour la porter à dix jours, comme
pour une transaction non compensée, et observer la hausse de la marge initiale d'un facteur $\sqrt 2$~; puis
faire tomber le spread de financement à zéro et expliquer pourquoi la base disparaît.
\end{tutorial}

\section{Construction~: le calculateur de marge et de financement}\label{bld:m2:swap-clearing-and-the-ccp-basis:mva}

\begin{build}
\bfield{Rôle} La petite société de bourse cote des swaps auprès de clients de différentes
chambres de compensation et doit savoir ce que coûte chaque transaction en marge dans chaque
chambre~; son système de risque (livre~6) étend ce calcul
aux modèles de marge réels.
\bfield{Interface} \texttt{MarginModel(sigma\_bp\_day, mpor\_days, z)} avec
\texttt{im\_per\_dv01()}~; \texttt{annuity\_remaining(t, maturity, rate)}~;
\texttt{dv01\_path}~; \texttt{im\_path}~; \texttt{mva(notional, maturity, rate,
model, funding\_spread)}~; \texttt{basis\_bp(\ldots, ccps)}.
\bfield{Règles} Courbe plate, périodes annuelles~; marge proportionnelle au DV01 du swap
restant~; pas trimestriels pour l'intégrale de financement~; la base par position
unidirectionnelle, multipliée par le nombre de chambres dans lesquelles le dealer est
unidirectionnel.
\bfield{Tests d'acceptation} \texttt{code/firm/ccpbasis/tests/}~: la forme fermée de
l'annuité~; marge proportionnelle au DV01 et décroissante~; MVA linéaire en notionnel et en
financement~; base croissante avec la maturité et additive entre les chambres.
\bfield{Pour aller plus loin} Une marge de portefeuille (compensation entre transactions dans
une même chambre)~; une marge par simulation historique sur la courbe de
la \cref{bld:m2:interest-rate-swaps:curve}~; la marge incrémentale d'une nouvelle transaction
dans chaque chambre, qui est ce que le dealer tarifie réellement.
\end{build}

\begin{omsources}
\item G20, \emph{Leaders' Statement: The Pittsburgh Summit}, septembre 2009.
\item Commodity Futures Trading Commission, communiqués de presse sur la première
détermination de compensation (novembre 2012) et sur ses phases (2013).
\item Comité de Bâle et OICV, \emph{Margin requirements for non-centrally
cleared derivatives}, et communiqués sur les phases finales (2019, 2020).
\item Clarus Financial Technology, ``CME-LCH basis for dummies'', juin 2017~; DLA
Piper, ``EMIR 3 active account requirement'', juin 2025.
\end{omsources}

\section{Exercices}

\begin{exercise}[$\star$]\label{exo:m2:swap-clearing-and-the-ccp-basis:1}
Dans le modèle de l'\cref{ex:m2:swap-clearing-and-the-ccp-basis:mva}, calculer la marge initiale
qu'exige un nouveau swap de dix ans d'un notionnel de 100~millions de dollars.
\end{exercise}

\begin{exercise}[$\star$]\label{exo:m2:swap-clearing-and-the-ccp-basis:2}
Un client reçoit le taux fixe de ce swap. Les taux montent de 3 points de base en une journée. Calculer
la marge de variation qu'il verse.
\end{exercise}

\begin{exercise}[$\star$]\label{exo:m2:swap-clearing-and-the-ccp-basis:3}
Le courtier compensateur d'un client fait défaut. Dire ce qu'il advient des swaps du client si le
portage réussit, et s'il échoue.
\end{exercise}

\begin{exercise}[$\star\star$]\label{exo:m2:swap-clearing-and-the-ccp-basis:4}
Calculer la MVA de 100~millions de dollars du swap de dix ans pour un spread de financement
de 50 points de base, ainsi que la base courante qu'elle implique.
\end{exercise}

\begin{exercise}[$\star\star$]\label{exo:m2:swap-clearing-and-the-ccp-basis:5}
Dans une chambre, les clients sont en majorité payeurs de taux fixe. Dans quelle chambre le \omterm{def:m2:interest-rate-swaps:par}{taux de swap
au pair} est-il plus élevé, et pourquoi~?
\end{exercise}

\begin{exercise}[$\star\star$]\label{exo:m2:swap-clearing-and-the-ccp-basis:6}
Un fonds détient 9~milliards d'euros de dérivés non compensés en cours et un modèle de marge qui
donne une marge initiale de 70~millions d'euros face à un dealer. Est-il soumis aux \omterm{def:m2:swap-clearing-and-the-ccp-basis:umr}{règles de marge
pour les dérivés non compensés}~? Quel montant de marge initiale est échangé~?
\end{exercise}

\begin{exercise}[$\star\star\star$]\label{exo:m2:swap-clearing-and-the-ccp-basis:7}
\emph{Codage.} À l'aide de \texttt{basis\_bp}, calculer la base à trente ans par position
unidirectionnelle pour des spreads de financement de 25, 50 et 75 points de base. Comparer le total
pour deux chambres à 50 avec les 3.4 points de base observés en 2017.
\end{exercise}

\begin{exercise}[$\star\star\star$]\label{exo:m2:swap-clearing-and-the-ccp-basis:8}
\emph{Repérer l'erreur.} « Le swap à trente ans est de 3 points de base plus cher dans une chambre
que dans l'autre~: recevoir dans la première, payer dans la seconde, et empocher 3 points de base
par an sans risque. » Corriger ce raisonnement.
\end{exercise}

\section{Problème~: le livre unidirectionnel}

\begin{problem}[{Problème du week-end --- ce qu'un dealer facture pour une chambre déséquilibrée}]\label{pb:m2:swap-clearing-and-the-ccp-basis:1}
Un dealer a reçu le taux fixe sur 2~milliards de dollars de swaps de dix ans auprès de clients de la
chambre A, où les clients ne reçoivent presque jamais. Il se couvre en payant le taux fixe sur
2~milliards de dollars dans la chambre B. Utiliser le modèle stylisé de
la \cref{prop:m2:swap-clearing-and-the-ccp-basis:mva}~: 7 points de base par jour, cinq
jours, $z = 2.326$, une courbe plate à 4~\%, un spread de financement de 50 points de base.

\medskip\noindent\textbf{Partie I --- Chambre A.}
\begin{enumerate}
\item Donner le DV01 de la position dans la chambre A.
\item Donner sa marge initiale.
\item Donner la marge de variation pour une hausse des taux de 5 points de base, et préciser qui la verse.
\item Pourquoi la chambre ne se soucie-t-elle pas du fait que le dealer est couvert ailleurs~?
\item Comment la marge changerait-elle si la position du dealer en A était compensée par
les receveurs d'autres dealers~?
\end{enumerate}

\medskip\noindent\textbf{Part II --- Both houses.}
\begin{enumerate}[resume]
\item Donner la marge initiale à la maison B.
\item Donner la MVA de chaque position, puis celle des deux positions réunies.
\item Donner la base courante qui rémunère une position, puis celle des deux.
\item Qu'en serait-il pour des swaps à trente ans~?
\item Quelle serait la marge initiale avec une période de marge de dix jours~?
\end{enumerate}

\medskip\noindent\textbf{Partie III --- Le choix.}
\begin{enumerate}[resume]
\item Le dealer pourrait à la place se couvrir à la maison A, si un autre dealer y recevait
le taux fixe. Quelle marge déposerait-il alors~?
\item Quel est le montant maximal que le dealer devrait verser à un autre dealer, en points de base
de taux, pour recevoir le taux fixe de sa part en A plutôt que de se couvrir en B~?
\item Pourquoi aucun dealer ne voudrait-il conclure cette transaction~?
\item Comment les clients finissent-ils par payer~?
\item Qu'est-ce qui modifierait le déséquilibre à la maison A~?
\end{enumerate}

\medskip\noindent\textbf{Partie IV --- Jugement.}
\begin{enumerate}[resume]
\item Pourquoi la marge déposée en espèces coûte-t-elle quelque chose~?
\item Pourquoi les marges des maisons réelles diffèrent-elles de ce modèle~?
\item Comment une règle exigeant un compte auprès d'une maison donnée déplace-t-elle la base~?
\item Énoncer le \emph{résultat nommé}~: la base à laquelle se couvrir à la maison B et compenser
à la maison A coûtent le même prix, pour des swaps à dix ans.
\item En une phrase~: pourquoi des swaps identiques ont-ils deux prix~?
\end{enumerate}
\end{problem}

\section{Questions d'entretien}

\begin{interviewq}[$\star$ \iqroles{banque, trader}]\label{iq:m2:swap-clearing-and-the-ccp-basis:1}
Pourquoi les swaps ont-ils été transférés à la compensation centrale après 2008, et en quoi
cela a-t-il changé pour un dealer~?
\end{interviewq}

\begin{interviewq}[$\star$ \iqroles{banque, développeur}]\label{iq:m2:swap-clearing-and-the-ccp-basis:2}
Quelle est la différence entre la marge initiale et la marge de variation~?
\end{interviewq}

\begin{interviewq}[$\star\star$ \iqroles{trader, banque}]\label{iq:m2:swap-clearing-and-the-ccp-basis:3}
Expliquer la base CME--LCH. S'agit-il d'un arbitrage~?
\end{interviewq}

\begin{interviewq}[$\star\star$ \iqroles{chercheur, banque}]\label{iq:m2:swap-clearing-and-the-ccp-basis:4}
Qu'est-ce que la MVA, et comment l'estimeriez-vous pour un nouveau swap~?
\end{interviewq}

\begin{interviewq}[$\star\star$ \iqroles{banque}]\label{iq:m2:swap-clearing-and-the-ccp-basis:5}
Un client négocie un swap de gré à gré sous CSA au lieu de le compenser. Quelles sont
les différences dans ses coûts et ses risques~?
\end{interviewq}

\begin{interviewq}[$\star\star\star$ \iqroles{développeur, chercheur}]\label{iq:m2:swap-clearing-and-the-ccp-basis:6}
Comment calculeriez-vous la marge initiale incrémentale d'une nouvelle transaction dans
une chambre de compensation, assez vite pour coter un client~?
\end{interviewq}
